\documentclass[10pt,a4paper]{amsart}
\usepackage[margin=19mm]{geometry}
\usepackage{amsmath,amssymb,mathtools}
\usepackage[scr=rsfs,cal=euler]{mathalfa}
\usepackage{xfrac}
\usepackage[unicode,hidelinks,bookmarks=false]{hyperref}
\newcommand{\ie}{i.e.\ }
\newcommand{\cf}{cf.\ }
\newcommand{\resp}{respectively\ }
\newcommand{\Subsection}{\subsection}
\providecommand{\nobreakdash}{\nobreak\hbox{-}}
\setlength{\emergencystretch}{2em}
\multlinegap0pt
\usepackage{amssymb}
\usepackage{tikz}
\usepackage{float}
\newtheorem{lemma}{Lemma}
\newtheorem{remark}{Remark}
\newtheorem{theorem}{Theorem}
\newcommand{\C}{\mathcal{C}}
\newcommand{\acl}{\mathrm{acl}}
\renewcommand{\d}{\delta}
\newcommand{\dcl}{\mathrm{dcl}}
\renewcommand{\epsilon}{\varepsilon}
\title{On the model theory of the Farey graph}
\author{Zahra Mohammadi Khangheshlaghi, Katrin Tent}
\date{Source: arXiv:2503.02121v2 (CC BY 4.0)}
\begin{document}
\maketitle

\noindent\emph{Source and licence.} On the model theory of the Farey graph. Version arXiv:2503.02121v2 (CC BY 4.0), distributed by arXiv under the Creative Commons Attribution 4.0 International licence, http://creativecommons.org/licenses/by/4.0/, as recorded in the arXiv licence metadata for this exact version and shown on the abstract page https://arxiv.org/abs/2503.02121v2. Full licence notice: \texttt{licenses/arxiv-2503.02121-CC-BY-4.0.txt}.\par\smallskip

\noindent\textit{Editorial scope.} Counted unit: the article's theorem thm:omega (source0.tex lines 588--590). The article's complete original proof of it is delivered with it (source0.tex lines 613--665, one \texttt{proof} environment of 53 lines, which the article places away from the statement). No mathematics is written, repaired or re-derived: the statement and the proof are the article's own lines, in the article's own notation, and no retained line is substituted or re-flowed.  Retained uncounted as the source-local context the proof needs: lines 364--368; lines 482--485; lines 516--520; lines 594--596; lines 602--604.  Nothing was omitted from any retained span, so the package carries no omission record.  No retained line contains a citation, so the wrapper carries no bibliography and no bibliography entry.  No retained line contains a figure environment, an image-inclusion command or drawing code, so no image is delivered and the package declares no asset dependency.  Editorial formatting only, in addition: an AMS \texttt{amsart} preamble that restates the article's own declarations and macros used by the retained lines, namely the environments \texttt{lemma}, \texttt{remark}, \texttt{theorem} on separate counters, together with the standard packages those lines need.  The retained environments print in this document as Lemma 1, Remark 1, Theorem 1 in order of appearance, whereas the article prints the same items with the numbering of its own full document.  The complete native source of the article as distributed in the arXiv e-print for this exact version is bundled unchanged as \texttt{sources/174174/source0.tex}.
\begin{lemma}\label{lem:finitely_many_path}
Let $M$ be a model of $T$.
    For  any two vertices $x, y\in M$ 
  there are only finitely many geodesics from $x$ to $y$.
\end{lemma}

\begin{remark}\label{rem:unique connecting points}
Let $M$ be a model of $T$ and suppose that $m$ is minimal such that for some $\delta=(C_1,\ldots, C_m)$ we have $P_\d(a,b)$. Then the connecting points $z_0=a,\ldots, z_m=b$ are uniquely determined and are the vertices of type $M$ in $G_{tree}(M)$ on the shortest path from $a$  to $b$, i.e. $z_0,\ldots, z_m\in\dcl(a,b)$.
If $m$ is not necessarily minimal, then by Lemma~\ref{lem:finitely_many_path} the connecting points are in $\acl(a,b)$
\end{remark}

\begin{remark}\label{rem:Y algebraic}
Note that if $\epsilon=(\d_1,\d_2,\d_3,n,\sigma)$ has minimal total length such that $Y_\epsilon(x,a,b)$ holds for some witnesses $x',a',b'$, then by the uniqueness from Remark~\ref{rem:unique connecting points} and the algebraicity of $D_{n,\sigma}(x',a',b')$  we see that $x'$ is the gate for $x$ to $\acl(a,b)$.
In particular, if $\d_1=\emptyset$, then $x\in\acl(a,b)$.
 
\end{remark}

\begin{lemma}\label{lem:finitely many cycles}
Let $M$ be a model of $T$ and let $(x,y)$ be an edge, $3\leq k<\omega$. There are finitely many simple $k$-cycles  containing $(x,y)$.
\end{lemma}

\begin{lemma}\label{lem:str min}
Let $a,b\in G_F, k=d(a,b),$ and $C\in \C_n, C'\in\C_m$. Then the formula $P_C(x,a)\wedge P_{C'}(x,b)$ is algebraic.
\end{lemma}

 \begin{theorem}\label{thm:omega}
     The Morley rank of the theory of the Farey graph is $\omega$.
 \end{theorem}

 \begin{proof}[Proof of Theorem~\ref{thm:omega}]


By quantifier elimination in the language $L'$, the definition of the predicates $Y_{\delta_1, \delta_2, \delta_3, n, \sigma}$ as a conjunction of instances of formulas $P_\delta$ and Remark~\ref{rem:Y algebraic} it is enough to prove that for any $a\in M$ the Morley rank of     $P_\delta(x,a)$ is exactly $m=|\delta|$. 

    Let $M$ be an $\omega$-saturated model of $T$, so $M$ has infinitely many connected components and every vertex in $G_{tree}(M)$ has infinite valency.

The proof is by 
 induction on $m$.  

\bigskip

Let $m=1$ and  let $F_i, i<\omega$, be  neighbours of $a$ in $G_{tree}(M)$. Since all $F_i$ are isomorphic to the Farey graph, each contains a vertex $b_i$ such that $P_\d(a, b_i)$ holds, hence  $RM(P_\d(a, x)) ) \geq 1$ for all $a\in M$.

We claim that $P_\d(x,a)$ is strongly minimal. It suffices to prove that for all sequences $\d'$ and all $b\in M$ the formula $P_{\d'}(x,b)\wedge P_\d(x,a)$ is either algebraic or cofinite in $P_\d(x,a)$.

So let $|\d'|=m'$ be minimal such that $P_\d(a,x)\wedge P_{\d'}(x,b)$ holds for infinitely many $x$. Then $m'\geq 2$ by Lemma~\ref{lem:str min}.
If $a, b$ lie in a different equivalence classes, then by Remark~\ref{rem:unique connecting points} the set of solutions for $P_\d(a,x)\wedge P_{\d'}(x,b)$ is the same as the set of solutions $P_\d(a,x)\wedge P_{\d''}(x,c)$ where $c$ is the gate for $b$ to the equivalence class of $a$ and $\d''$ is a proper end segment of $\d'$. This contradicts the minimality of $m'$. So $a, b$ lie in the same Farey graph.

If $\d'=\d_1\circ\d_2$, $P_{\d_1}(a,b)$ and $\d_2$ embeds into $\d$, then clearly $P_\d(x,a)\wedge P_{\d'}(x,b)$ is confinite $P_\d(x,a)$. 

If $P_{\d_1}(a,b)$ and $\d_2$ does not embed into $\d$, then $P_{\d_2}(x,a)\wedge P_\d(x,a)$ is finite by minimality of $m'$. 

Hence we now assume we cannot write $\d'=\d_1\circ\d_2$ such that $P_{\d_1}(a,b)$ holds.

Consider the finitely many geodesics from $a$ to $b$ and write $\d'=(C)\circ\d''$. 
There are at most  finitely many copies of $\d$ and of $C$ intersecting one of these geodesics by Lemma~\ref{lem:finitely many cycles}. 
If all copies of $C$ intersect one of these geodesics, we could replace $P_{\d'}(x,b)$ by $P_{\d''}(x,c)$ where $c$ is the corresponding connecting point, contradicting minimality of $m'$.
Hence infinitely many copies of $C$ do not intersect any of these geodesics.
Any realization of $P_\d(a,x)\wedge P_{\d'}(x,b)$ arises from a unique copy of $\d$ and a unique copy of $\d'$. Since these copies form a simple cycle of bounded length containing a geodesic from $a$ to $b$, there are only finitely many copies of $C$ involved. Hence we can again replace $P_{\d'}(x,b)$ by $P_{\d''}(x,c)$ where $c$ is the corresponding connecting point, contradicting minimality of $m'$.

Now suppose the claim is proven for all $m'<m$ and suppose $\d=(C_1,C_2)\circ\d'$ where $|\d'|=m-2$. Let $b_i\in F_i$ be such that $P_C(a,b_i)$ holds. Then by induction assumption, $P_{C\circ \d'}(x,b_i)$ has Morley rank at least $m-1$. Since the $b_i$ lie in distinct copies of the Farey graph, and any path from $b_i$ to $b_j$ passes through $a$, the sets $P_{\d'}(x,b_i)$ and $P_{\d'}(x,b_j)$ intersect at most in the set $P_{\d'}(x,a)$. Again by induction assumption these sets have Morley rank $m-2$, showing that $P_\d(x,a)$ has Morley rank at least $m$.


We now show that for any $\C$-sequence $\d'$ and vertex $b$, the Morley rank of either $P_{\d'}(x,b)\wedge P_\d(x,a)$ or  $\neg P_{\d'}(x,b)\wedge P_\d(x,a)$ is strictly less than $m$. 

 
Let $m'$ be minimal such that for some $\d'$ of length $m'$  this is not the case.

If $a,b$ lie in different equivalence classes, the result follows from the uniqueness in Remark~\ref{rem:unique connecting points} and the induction assumption.



If $\d'=\d_1\circ\d_2$, $P_{\d_1}(a,b)$ and $\d_2$ embeds into $\d$, then  $\neg P_{\d'}(x,b)\wedge P_\d(x,a)$ has Morley rank strictly less than $m$.



If $P_{\d_1}(a,b)$ and $\d_2$ does not embed into $\d$, then $P_{\d_2}(x,a)\wedge P_\d(x,a)$ has smaller Morley rank by minimality of $m'$. 

Hence we now assume we cannot write $\d'=\d_1\circ\d_2$ such that $P_{\d_1}(a,b)$ holds.

Consider the finitely many geodesics from $a$ to $b$ and write $\d'=(C)\circ\d''$. Now we argue exactly as in the case $m=1$. This concludes the proof.
\end{proof}

\end{document}
